\chapter{Calculer avec les décimaux}\label{ch:g5:decimalops}

Les nombres décimaux s'additionnent, se soustraient et se
multiplient presque comme les entiers --- tout l'art est de savoir
où va la \omterm{def:g4:decimals:point}{virgule}. Ce chapitre couvre les méthodes en colonnes et la
magie de la multiplication ou de la division par $10$, $100$,
$1000$. (Multiplier deux décimaux entre eux attend
\cref{ch:g6:decimals}.)

\section{Additionner et soustraire}

\begin{method}[Colonnes avec une virgule]\label{met:g5:decimalops:addsub}
\begin{enumerate}
  \item Écrire les nombres avec les \emph{\omterm{def:g4:decimals:point}{virgules} alignées}
        (alors les unités sont sous les unités, les dixièmes sous
        les dixièmes)~;
  \item compléter la partie décimale plus courte avec des \omterm{def:g1:counting:zero}{zéros}~;
  \item additionner ou soustraire comme avec les entiers~;
  \item faire descendre la \omterm{def:g4:decimals:point}{virgule} du résultat droit en bas.
\end{enumerate}
\end{method}

\begin{example}\label{ex:g5:decimalops:addsub}
$27.5 + 3.86$ (compléter $27.5 = 27.50$)~:
\[
\begin{array}{r}
2\,7.5\,0 \\
+\ \ \ 3.8\,6 \\
\hline
3\,1.3\,6
\end{array}
\qquad\qquad
\begin{array}{r}
1\,2.4\,0 \\
-\ \ \ 5.6\,5 \\
\hline
\ \ \,6.7\,5
\end{array}
\]
Pour la soustraction ($12.4 - 5.65$)~: compléter, emprunter comme
d'habitude, et vérifier en additionnant~: $6.75 + 5.65 = 12.4$.
\end{example}

\begin{example}[Compléments mentaux]\label{ex:g5:decimalops:complement}
Que faut-il ajouter à $7.6$ pour atteindre $10$~? Monter~:
$7.6 \to 8$ fait $0.4$, puis $8 \to 10$ fait $2$~: le complément
est $2.4$. Pratique pour la monnaie~: payer $10$ pour un article à
$7.60$, recevoir $2.40$ de rendu.
\end{example}

\section{Dix fois plus grand, dix fois plus petit}

\begin{proposition}[Fois et divisé par 10, 100, 1000]\label{prop:g5:decimalops:shift}
Multiplier par $10$ fait valoir chaque chiffre dix fois plus~: les
chiffres glissent d'une place à gauche --- ce qui ressemble à la
\emph{\omterm{def:g4:decimals:point}{virgule} qui avance d'une place à droite}. Diviser fait
l'inverse~:
\[
3.47 \times 10 = 34.7, \qquad
3.47 \times 100 = 347, \qquad
52.1 \div 10 = 5.21, \qquad
6 \div 100 = 0.06 .
\]
\end{proposition}
\admitted

\begin{omfigure}
\begin{tikzpicture}[scale=0.9]
  \draw (0,0) grid[xstep=1.5, ystep=0.75] (6.0,2.25);
  \node[font=\small] at (0.75,2.6) {dizaines};
  \node[font=\small] at (2.25,2.6) {unités};
  \node[font=\small] at (3.75,2.6) {dixièmes};
  \node[font=\small] at (5.25,2.6) {centièmes};
  \node[font=\normalsize, omExo] at (2.25,1.85) {$3$};
  \node[font=\normalsize, omExo] at (3.75,1.85) {$4$};
  \node[font=\normalsize, omExo] at (5.25,1.85) {$7$};
  \node[font=\normalsize, omDef] at (0.75,1.1) {$3$};
  \node[font=\normalsize, omDef] at (2.25,1.1) {$4$};
  \node[font=\normalsize, omDef] at (3.75,1.1) {$7$};
  \node[right, font=\small] at (6.2,1.85) {$3.47$};
  \node[right, font=\small, omDef] at (6.2,1.1)
    {$3.47 \times 10 = 34.7$};
  \draw[-Stealth, omDef, thick] (4.4,1.7) -- (3.1,1.25);
\end{tikzpicture}

{\small Multiplier par $10$~: chaque chiffre monte d'une place à
gauche. La \omterm{def:g4:decimals:point}{virgule} ne bouge pas vraiment --- ce sont les chiffres
qui bougent.}
\end{omfigure}

\begin{example}[Les conversions, encore]\label{ex:g5:decimalops:convert}
Les unités de mesure sont des puissances de dix les unes des
autres, donc les conversions sont juste ces décalages~: $2.35$~m
$= 235$~cm ($\times 100$)~; $470$~g $= 0.47$~kg ($\div 1000$)~;
$1.5$~L $= 150$~cL.
\end{example}

\section{Multiplier un décimal par un entier}

\begin{method}[Décimal fois entier]\label{met:g5:decimalops:mult}
Calculer comme s'il n'y avait pas de \omterm{def:g4:decimals:point}{virgule}, puis donner au
résultat autant de chiffres décimaux que le facteur décimal~:
\[
4.35 \times 6: \quad 435 \times 6 = 2\,610
\ \to\ 4.35 \times 6 = 26.10 = 26.1 .
\]
Estimer pour placer la \omterm{def:g4:decimals:point}{virgule} avec confiance~: $4.35$ est proche
de $4$, et $4 \times 6 = 24$ --- donc $26.1$, ni $2.61$ ni $261$.
\end{method}

\begin{example}[L'addition répétée marche toujours]\label{ex:g5:decimalops:mult}
$0.75 \times 4 = 0.75 + 0.75 + 0.75 + 0.75 = 3$~: quatre \omterm{def:g3:division:half}{quarts}
font un tout, en habits décimaux. Multiplier par un entier, c'est
toujours «~autant de fois~».
\end{example}

\begin{example}[Courses]\label{ex:g5:decimalops:shopping}
Trois cahiers à $2.45$ chacun et un stylo à $1.20$~:
\begin{enumerate}
  \item cahiers~: $2.45 \times 3 = 7.35$~;
  \item total~: $7.35 + 1.20 = 8.55$~;
  \item rendu sur $10$~: $10 - 8.55 = 1.45$.
\end{enumerate}
\end{example}

\section{Exercices}

\begin{exercise}[$\star$]\label{exo:g5:decimalops:1}
Calculer en colonnes~: $45.7 + 8.93$~; \quad $16.25 + 3.75$~; \quad
$0.86 + 12.4$.
\end{exercise}

\begin{exercise}[$\star$]\label{exo:g5:decimalops:2}
Calculer en colonnes et vérifier~: $9.4 - 2.72$~; \quad
$20 - 13.45$~; \quad $6.03 - 5.9$.
\end{exercise}

\begin{exercise}[$\star$]\label{exo:g5:decimalops:3}
Compléments mentaux à $10$ (comme dans
\cref{ex:g5:decimalops:complement})~: $6.2$~; \quad $3.75$~; \quad
$9.99$.
\end{exercise}

\begin{exercise}[$\star$]\label{exo:g5:decimalops:4}
Calculer sans colonnes~:
\[
7.24 \times 10, \qquad
0.58 \times 100, \qquad
34.5 \div 10, \qquad
7 \div 100, \qquad
0.3 \times 1000 .
\]
\end{exercise}

\begin{exercise}[$\star$]\label{exo:g5:decimalops:5}
Convertir~: $4.05$~m en cm~; \quad $325$~cm en m~; \quad $0.6$~kg
en g~; \quad $85$~cL en L.
\end{exercise}

\begin{exercise}[$\star$]\label{exo:g5:decimalops:6}
Estimer d'abord, puis calculer~: $6.2 \times 4$~; \quad
$3.45 \times 8$~; \quad $12.5 \times 6$.
\end{exercise}

\begin{exercise}[$\star$]\label{exo:g5:decimalops:7}
Un tour de piste mesure $0.4$~km. Quelle longueur pour $7$ tours~?
Et pour $25$ tours~?
\end{exercise}

\begin{exercise}[$\star$]\label{exo:g5:decimalops:8}
Léa achète $2$ baguettes à $1.15$ chacune et un gâteau à $12.60$.
Elle paie avec un billet de $20$. Calculer son total et sa monnaie.
\end{exercise}

\begin{exercise}[$\star$]\label{exo:g5:decimalops:9}
Une bouteille contient $1.5$~L. Combien de litres dans un pack de
$6$ bouteilles~? Une famille boit $0.75$~L par jour~: pendant
combien de jours un pack suffit-il~? (Combien de fois $0.75$
rentre-t-il dans $9$~?)
\end{exercise}

\begin{exercise}[$\star\star$]\label{exo:g5:decimalops:10}
Milo calcule $5.6 \times 3 = 15.18$, en raisonnant «~$5 \times 3 = 15$
et $6 \times 3 = 18$~». Utiliser une estimation pour montrer que la
réponse est forcément fausse, puis calculer correctement.
\end{exercise}

\begin{exercise}[$\star\star$]\label{exo:g5:decimalops:11}
Les quatre tours d'une coureuse ont pris $65.4$~s, $64.9$~s,
$65.0$~s et $66.1$~s.
\begin{enumerate}
  \item Calculer le temps total des quatre tours.
  \item Le record pour quatre tours est $260$~s. L'a-t-elle battu~?
        De combien~?
\end{enumerate}
\end{exercise}
